% Modernised reading — fonds Alexandre Grothendieck, shelfmark 158,
% pages 1-85 (the whole folder, its five batches read together in one pass).
%
% This is an INTERPRETATION, not a transcription. It re-reads the pages in
% today's notation and vocabulary, states the hypotheses the manuscript leaves
% implicit, and drops the critical apparatus so that the mathematics reads
% straight through. Everything the brackets and underlines carried — a doubtful
% reading, a notation he used differently, a missing passage — moves into a
% footnote.
%
% It is held to being mathematically correct as it stands, not merely faithful
% to the page. Where the manuscript is loose or mistaken, the true statement is
% what appears here and the footnote says what the page has. In any dispute of
% substance the transcription governs: whoever cites Grothendieck must cite
% batch-01.fr.tex through batch-05.fr.tex.
%
% Pass: Opus 5 (claude-opus-5), 2026-09-21 — first pass, unchecked against the
% pages by a human. Opus 5 rather than the skill's Fable 5.1 default, for
% continuity with folders 159 and 160 of the same inventory group and with the
% folder's own five transcriptions, all of which ran on Opus 5.
%
% Scope: a SINGLE document for the whole shelfmark. No part of folder 158 is
% read in another file.
%
% NOTATION FIXED FOR THE WHOLE DOCUMENT, and where it departs from the pages:
%   - X^ (presheaves on X) is written \widehat{X} throughout. The pages write a
%     raised circumflex, sometimes superimposed on the letter. His X^v, the
%     COVARIANT functors to sets, is written \widecheck{X}-style as X^{\vee}
%     only where the distinction is in play (page 61).
%   - X_{//a} is the full subcategory of objects x with Hom(a,x) non-empty —
%     the cosieve generated by a. The pages write a two-stroke double slash,
%     prefixed (beta//Z, page 62) or subscripted (X-tilde_{//xi-tilde},
%     pages 76 and 79); one notation is used here for both, and the reading
%     X_{//a} is fixed by page 57, where X_{//a} inter X_{//b} non-empty is
%     shown to be the equivalence relation defining pi_0. X_{//(a,b)} is the
%     intersection.
%   - x\X and X_{/x} for the comma categories under and over x. The pages put
%     the index to the LEFT of a backslash for the first; page 74 carries both
%     senses in one formula, which is what fixes the convention.
%   - M_0, G_0 = pi_1(M_0), delta, delta_0 as the pages have them.
%   - Delta-tilde is the category of non-empty finite sets (pages 6-15). The
%     plain Delta of page 17 is a GROUP and is unrelated; the collision of
%     letter is his, and the group is written Delta(E) here as there.
%   - W is a fundamental localiser; the pages number its axioms Loc(6),
%     Loc(6a), Loc(8) without stating them, and this reading does not invent
%     statements for them.
%
% SUBSTANTIVE DEPARTURES FROM THE PAGES, each footnoted where it occurs:
%   - page 22: the quantifier is drawn "for all" where the sense requires
%     "there exist"; the definition of M_0 is stated here with the correct one.
%   - page 23: he calls delta "multiplicative" and writes the additive law;
%     it is a homomorphism to the ADDITIVE group Z, and is named so here.
%   - page 28: the last factor of the boxed relation carries an exponent -1
%     that his own composition rule does not produce, and its opening
%     parenthesis is never closed. The relation is stated here as the rule
%     gives it.
%   - page 41: "the set Ep(E) of epimorphisms of E = the set of minimal
%     elements of E" defines E by itself; the argument requires a separate E_0,
%     which the page supplies under that name only at page 45.
%   - page 45: u(E_0) = E_0 needs E_0 FINITE, which the page uses without
%     saying so at that point.
%   - page 49: the Corollaire is repudiated in his own margin — "Faux ! revoir
%     la demonstration" — and is given here with that status, not as a result.
%   - page 52: the open subcategory of least objects is written X_0 in the
%     proposition and X^0 in the corollary. One notation is used here.
%   - page 78: the last sentence's subject should be x-tilde, not xi-tilde.
%   - page 81: the two underbrace sources are drawn differently (an m and an
%     a) where the argument wants one source; page 74 writes the same sum with
%     one source, and that is what is used here.
%
% WHAT THE FOLDER ANNOUNCES AND NEVER ESTABLISHES — recorded because an
% edition that reads continuously must not step over an absence:
%   - the Lemme (?) of page 10, explicitly conjectural, on which Corollaires 1
%     and 2 of that page depend;
%   - Proposition 2 (??) of page 46, whose proof is attempted on pages 47-48
%     and abandoned;
%   - Proposition 3 of page 35, whose statement breaks off before its
%     conclusion;
%   - the Theoreme of page 39, which breaks off inside its case a);
%   - the folder's closing question (pages 81-83), which is left open in his
%     own words.
\documentclass[11pt,a4paper]{article}
% Le préambule commun à toutes les transcriptions du fonds Grothendieck.
%
% Il tient en une idée : **l'incertitude se compose, elle ne se résout pas.**
% Une transcription qui lisse un mot douteux a détruit la seule chose pour
% laquelle on la faisait — un lecteur doit pouvoir savoir, à chaque endroit,
% ce qui était sur le feuillet et ce qui a été ajouté. Les six macros qui
% suivent sont donc l'appareil critique, et non de la décoration.
%
% Les mêmes commandes sont reconnues par `scripts/render.mjs`, qui produit la
% vue de lecture du site. Une macro ajoutée ici doit l'être aussi là-bas,
% faute de quoi le rendu échouera bruyamment — ce qui est voulu : un rendu
% silencieusement faux est pire qu'un rendu absent.

\ProvidesPackage{grothendieck}[2026/08/08 Transcriptions du fonds Grothendieck]

\usepackage{amsmath,amssymb,amsthm}
\usepackage{mathrsfs}
\usepackage{stmaryrd}
% Les diagrammes commutatifs sont partout dans ce fonds : c'est la langue
% même de la géométrie algébrique de Grothendieck, et une transcription qui
% les rendrait en prose ne transcrirait rien.
\usepackage{tikz-cd}
% Français par défaut — c'est la langue des feuillets — mais l'anglais est
% chargé aussi : la traduction et le résumé sont des documents anglais, et la
% typographie française (espaces avant les deux-points) y serait une faute.
% Ces documents basculent par \selectlanguage{english} en tête de corps.
\usepackage[english,french]{babel}
\usepackage{fontspec}
\usepackage{geometry}
\geometry{a4paper,margin=25mm,marginparwidth=28mm}
\usepackage{marginnote}
\usepackage{xcolor}
% ulem, not soul. `soul`'s \st and \ul are built on a character-by-character
% scan that XeLaTeX's font machinery breaks: the document fails with an
% undefined control sequence rather than degrading. ulem does the same two jobs
% and is Unicode-engine safe. `normalem` stops it hijacking \emph, which the
% transcriptions use for Grothendieck's own emphasis.
\usepackage[normalem]{ulem}
\usepackage{hyperref}
\hypersetup{colorlinks=true,linkcolor=black,urlcolor=black,pdfborder={0 0 0}}

% --- Métadonnées du lot -----------------------------------------------------
% Reprises telles quelles par le script de rendu, qui les affiche en tête de
% la vue de lecture. Elles fixent ce que le fichier prétend couvrir : sans
% elles, une transcription flotte sans point d'attache dans un fonds de seize
% mille feuillets.

\newcommand{\folder}[1]{\gdef\grothFolder{#1}}
\newcommand{\batch}[1]{\gdef\grothBatch{#1}}
\newcommand{\leaves}[2]{\gdef\grothFirst{#1}\gdef\grothLast{#2}}
\newcommand{\foldertitle}[1]{\gdef\grothFoldertitle{#1}}
\gdef\grothFolder{?}\gdef\grothBatch{?}\gdef\grothFirst{?}\gdef\grothLast{?}\gdef\grothFoldertitle{}

% --- Le résumé --------------------------------------------------------------
%
% Il ouvre la lecture modernisée, sous le titre « Résumé », et il est composé
% en petit. Ce n'est pas un ornement : c'est le seul passage du document qui
% ne soit pas des mathématiques, et un lecteur doit pouvoir le reconnaître
% avant de l'avoir lu — puis le sauter s'il connaît déjà le sujet, ou s'y
% arrêter s'il ne le connaît pas. Le filet qui le suit marque la reprise du
% corps.

\newenvironment{resume}
  {\small\setlength{\parskip}{0.45em}}
  {\par\medskip\hrule\medskip}

% --- Mots-clés ---------------------------------------------------------------
%
% En anglais, en clôture du résumé de la lecture modernisée : le vocabulaire
% moderne sous lequel chercher ce que les feuillets construisent. Le manifeste
% du site (`npm run manifest`) les extrait pour étiqueter la cote entière —
% cette ligne est la source unique des tags, il n'y a pas de fichier à côté.
\newcommand{\keywords}[1]{\par\smallskip\noindent
  {\footnotesize\textsc{Keywords} --- \textit{#1}}\par}

% --- L'appareil critique ----------------------------------------------------

% Le numéro de feuillet, en marge. C'est l'ancre qui permet de régler un
% désaccord sur une formule en regardant *un* feuillet plutôt qu'une chemise
% de six cents. Le site s'en sert aussi pour tourner les pages du fac-similé
% quand on fait défiler la transcription.
\newcommand{\leaf}[1]{%
  \marginnote{\footnotesize\color{gray!70}\textsf{#1}}%
  \label{leaf:#1}\ignorespaces}

% Illisible : on ne devine pas. Un mot inventé qui se lit comme les autres est
% le pire résultat possible.
\newcommand{\ill}{\textcolor{red!60!black}{[\,\ldots\,]}}

% Lecture douteuse : le mot est donné, et signalé comme incertain.
\newcommand{\uncertain}[1]{\uline{#1}}

% Ajout de l'éditeur — une lettre manquante, un mot sous-entendu.
\newcommand{\add}[1]{\textcolor{blue!50!black}{[#1]}}

% Biffé par Grothendieck lui-même. Ce qu'il a rayé fait partie du document :
% une rature dit souvent où la pensée a changé de direction.
\newcommand{\struck}[1]{\sout{#1}}

% Note du transcripteur, sur l'état matériel du feuillet.
\newcommand{\note}[1]{\par\smallskip{\small\color{gray!60!black}\textit{[#1]}}\par\smallskip}

% Note marginale de Grothendieck, distincte du corps du texte.
\newcommand{\marginal}[1]{\par\smallskip\noindent\hspace{1em}%
  {\small\color{gray!60!black}#1}\par\smallskip}

% --- Titre ------------------------------------------------------------------

\AtBeginDocument{%
  \begin{center}
    {\large\bfseries Fonds Alexandre Grothendieck --- cote n\textsuperscript{o}~\grothFolder}\\[2pt]
    {\itshape\grothFoldertitle}\\[4pt]
    {\small feuillets \grothFirst--\grothLast\ (lot \grothBatch)}\\[2pt]
    {\footnotesize\color{gray!60!black}Transcription établie d'après le fac-similé numérisé
     par l'Université de Montpellier.\\
     Les lectures incertaines sont soulignées, les illisibles notées [\,\ldots\,],
     les ajouts entre crochets.}
  \end{center}
  \vspace{1em}}

\endinput


\folder{158}
\pages{1}{85}
\dating{[vers 1990-1991]}
\watermark{Édition de démonstration\\interprétation personnelle de l'œuvre}
\foldertitle{[Autour des Dérivateurs] : notes manuscrites (s.d.) — lecture modernisée du dossier entier}

\begin{document}

\begin{resume}
Une catégorie, c'est un ensemble de choses et, entre elles, des façons
d'aller de l'une à l'autre que l'on sait composer. On peut la regarder comme
un réseau de routes entre des villes : partir de $a$, arriver en $b$, puis
continuer vers $c$. Deux propriétés d'un tel réseau vont s'opposer dans tout
ce dossier, et ce sont les deux seules qu'il faut avoir en tête.

La première est une propriété d'ordre, presque de bon sens : dit-on que, de
deux villes quelconques, on peut toujours rejoindre une même ville plus loin ?
Un réseau où c'est le cas est dit \emph{cofiltrant}. C'est la condition qui
permet de passer à la limite, de réconcilier deux descriptions partielles en
les poussant assez loin. C'est aussi, très exactement, la condition sous
laquelle une catégorie se comporte comme un système de mesures de plus en plus
fines d'une même chose.

La seconde est une propriété de forme. À toute catégorie on sait associer un
espace — on recolle un point par ville, un segment par route, un triangle par
composition — et cet espace a une forme, avec ou sans trous. « Sans trou de
dimension $0$ » veut dire d'un seul tenant ; « sans trou de dimension $1$ »
veut dire que tout lacet peut être rempli. On dit alors $0$-connexe et
$1$-connexe.

Le dossier pose une question, et une seule : \textbf{l'absence de trous
force-t-elle la convergence ?} Si un réseau n'a pas de lacet que l'on ne
puisse remplir, est-il forcément cofiltrant ? La question n'a rien d'oiseux.
Elle dit si une condition de forme, qui se teste par des invariants
calculables, suffit à garantir une condition d'ordre, qui est celle dont on a
besoin pour travailler.

La réponse, telle que ces feuillets la dégagent, est \emph{oui dans le cas
fini, non en général}, et le dossier consiste presque entièrement à établir
ces deux moitiés. Pour la première il descend au cas le plus simple — une
catégorie à un seul objet, c'est-à-dire un monoïde — et montre qu'un monoïde
fini, dans lequel deux éléments se réconcilient toujours plus loin et dont le
groupe engendré est trivial, possède un élément $p$ que tout absorbe
($up = p$ pour tout $u$) : la convergence est alors immédiate. Pour la
seconde il fabrique un contre-exemple, et c'est le morceau de bravoure du
dossier : partant d'un monoïde et de deux éléments $\varphi$ et $\psi$ bien
choisis, il construit une catégorie d'un seul tenant dont le groupe
fondamental se surjecte sur $\mathbf{Z}$ — un lacet que rien ne remplit. Les
monoïdes qui font marcher la machine sont les surjections d'un ensemble
infini, et les injections : deux monoïdes où l'on peut toujours se ramener
plus loin, jamais absorber.

Une troisième moitié, si l'on peut dire, occupe la fin du dossier et en est
la partie la plus intéressante. Il y reprend la question par la géométrie.
Partant d'un tout petit diagramme ordonné — deux chemins parallèles de $a$ à
$b$ — il fabrique une catégorie $X'$ et montre qu'elle est une
\emph{suspension itérée} : l'opération qui, d'un cercle, fait une sphère, et
d'une sphère la suivante. Tout se joue alors sur un seul objet $T$, la partie
du réseau que $\tilde a$ et $\tilde b$ atteignent tous deux. Si $T$ est vide,
$X'$ a la forme de la sphère $S^{2}$, et cette sphère est un témoin : elle
distingue entre elles les notions d'équivalence que l'on peut raisonnablement
imposer aux catégories. Le dossier s'arrête sur la tentative de décider si
$T$ peut être vide, et sur l'aveu qu'il ne voit pas comment conclure.

Ce qu'il faut chercher ensuite sous des noms modernes : les
\emph{localisateurs fondamentaux}, qui sont justement les notions
d'équivalence en question et que le dossier manipule sans les définir ; la
théorie des \emph{catégories test} et des \emph{dérivateurs}, dont ces
feuillets sont contemporains ; et, du côté des monoïdes, le
\emph{groupe enveloppant} et le \emph{noyau de Rees} d'un monoïde fini.

Un mot enfin sur ce que ces pages laissent voir de leur auteur au travail,
car c'est ce qu'elles ont de moins remplaçable. On l'y voit se donner une
généralité, la tenir quelques feuillets, puis y renoncer en toutes lettres :
« Je crois que je vais craquer, et restreindre la généralité ». On l'y voit
avouer une difficulté : « j'ai un peu mal à construire un $F$ $0$-connexe qui
ne soit $1$-connexe ! ». On l'y voit surtout se corriger : au bas d'un
corollaire soigneusement démontré, il revient écrire dans la marge
« Faux ! revoir la démonstration », et recommencer quarante lignes plus loin
la même proposition sur d'autres bases. Aucun des trois moments ne serait
lisible dans un texte publié, et les trois disent la même chose : le travail
consiste à se tromper d'une manière qu'on puisse rattraper.

\keywords{fundamental localiser, Grothendieck derivator, test category, aspheric functor, totally aspheric category, cofiltered category, pseudo-filtered category, simply connected category, classifying space of a monoid, group completion, Rees kernel, left zero semigroup, idempotent, surjective endomorphism monoid, injective endomorphism monoid, Eilenberg swindle, suspension of a category, comma category, cosieve, local system on a category, presheaf topos, induced topos, aspheric morphism of topoi, connected component functor, semi-simplicial object, free group presentation, fundamental group of a category, Quillen Theorem A, Thomason model structure, nerve of a category, fibred category, cofibred category, cohomology with local coefficients}
\end{resume}

\pagerange{1}{85}
\section*{Le fil du dossier, et les conventions}

Le dossier n'est pas un texte suivi et ne porte aucun titre de sa main. C'est
une campagne de travail, menée sur quelques semaines, autour d'une question
unique, et reprise sept fois par des chemins différents. Les feuillets ne se
suivent pas toujours : il y a des essais, des moutures abandonnées, et deux
feuillets reliés dans l'ordre inverse du sien.

\pagerange{1}{85}
\subsection*{La question}

Soit $X$ une petite catégorie et $\widehat{X}$ la catégorie des préfaisceaux
sur $X$. Le dossier tourne tout entier autour de l'implication

\begin{quote}
\emph{si tout objet $0$-connexe de $\widehat{X}$ est $1$-connexe, alors $X$
est cofiltrante,}
\end{quote}

de sa réciproque, et des invariants qui permettent de la tester. Les deux
notions sont celles du résumé : $X$ est \emph{cofiltrante} si deux objets
quelconques admettent un objet vers lequel ils s'envoient tous deux et si deux
flèches parallèles s'égalisent plus loin ; elle est \emph{pseudo-cofiltrante}
si seule la première clause est demandée.\footnote{Le dossier emploie
« pseudo-cofiltrant » sans le définir, et donne page 55 l'équivalent qui fixe
le sens : $\forall u, v$, $\exists u', v'$ avec $uu' = vv'$. Page 19 il note
l'équivalence avec « $\widehat{X}$ localement irréductible », qui est la
forme sous laquelle la condition se lit dans le topos.} Toute la difficulté
tient dans l'écart entre les deux : se ramener plus loin n'est pas absorber.

\pagerange{1}{85}
\subsection*{Les stations}

\begin{itemize}
\item \textbf{Pages 1 à 4} — feuillets d'essai. Construction de la catégorie
$\widetilde{\pi}\backslash B$ attachée à un foncteur $\pi$, et du foncteur
canonique qu'elle envoie dans un topos de préfaisceaux. C'est l'outil, monté
avant qu'on sache à quoi il servira.
\item \textbf{Pages 6 à 15} — sa pagination $1$ à $9$. L'hypothèse
« $0$-connexe entraîne $1$-connexe » est posée pour la première fois, et l'on
en tire un morphisme de topos $f$ dont on montre qu'il est
$W_{0}$- puis $W_{1}$-asphérique. La conclusion visée est un lemme qu'il
énonce avec un point d'interrogation et ne démontre pas.
\item \textbf{Pages 17 à 20} — sa pagination $0$, $1$, $2$. Une présentation
par générateurs et relations du groupe fondamental d'une catégorie. C'est
l'outil de calcul dont la suite se servira sans cesse.
\item \textbf{Pages 21 à 41} — sa pagination $3$ à $18$, puis $7$. Le
contre-exemple : un monoïde, deux éléments $\varphi$ et $\psi$, un degré à
valeurs dans $\mathbf{Z}$, et une catégorie $0$-connexe non $1$-connexe. C'est
le massif le plus long et le plus achevé du dossier.
\item \textbf{Pages 44 à 53} — sa pagination $1$ à $8$. Le cas fini, dans les
deux sens : un théorème sur les monoïdes finis, puis sa transposition aux
catégories finies, corrigée en cours de route et refaite sur d'autres bases.
\item \textbf{Pages 55 à 58} — sa pagination $1$ à $3$. Les deux exemples
infinis, $\mathrm{Ep}(E)$ et $\mathrm{Mon}(E)$, qui font tourner le
contre-exemple.
\item \textbf{Pages 60 à 79} — sa pagination $1$ à $11$. Les localisateurs
fondamentaux, l'asphéricité, et le calcul de suspensions qui ramène tout à un
seul objet $T$.
\item \textbf{Pages 81 et 83} — sa pagination $12$ et $13$. La tentative
d'interpréter $T$, et l'arrêt du dossier.
\end{itemize}

Deux avertissements sur cette liste. Sa pagination recommence cinq fois, et
deux suites portent des numéros qui se recouvrent : le feuillet 41 porte
« $7$ » alors qu'il achève, pour le sens, la suite numérotée $3$ à $18$ des
feuillets 21 à 39.\footnote{Le désaccord est réel et n'est pas résolu ici. La
transcription le signale au feuillet 41. Le raccord du sens, lui, ne fait pas
de doute : le feuillet 39 s'arrête sur « a) Cette opération fait sans points
fixes » et le feuillet 41 reprend exactement le calcul de ce cas, puis donne
le cas b).} Et le dossier restart deux fois sur lui-même : les feuillets 50 à
53 refont la proposition des feuillets 44 à 46 sur d'autres hypothèses, et le
feuillet 77 refait avec un diagramme plus simple l'argument du feuillet 74.
Ce ne sont pas des répétitions à fondre, et cette lecture les tient séparées.

\pagerange{1}{85}
\subsection*{Les conventions}

$\widehat{X}$ est la catégorie des préfaisceaux sur $X$, $X^{\vee}$ celle des
foncteurs covariants à valeurs ensembles ; la distinction ne sert qu'au
feuillet 61 et elle y est essentielle.\footnote{Il écrit les deux d'un accent
porté sur la lettre, et rien sur le feuillet ne les distingue que la glose. La
variance est ici ce qui décide du sens de l'énoncé, et elle est rétablie
d'après cette glose.} Pour $a$ un objet de $X$, on écrit $X_{/\!/a}$ la
sous-catégorie pleine des objets $x$ tels que $\mathrm{Hom}(a,x)$ soit non
vide, et $X_{/\!/(a,b)} = X_{/\!/a} \cap X_{/\!/b}$.\footnote{Le dossier
trace un signe à deux jambages obliques, tantôt en préfixe, tantôt en indice ;
une seule notation est retenue ici. La lecture est fixée par le feuillet 57,
où il montre que la relation « $X_{/\!/a} \cap X_{/\!/b} \neq \emptyset$ » est
l'équivalence définissant $\pi_{0}(X)$ — ce qui n'a de sens que pour le
cocrible engendré.} Enfin ${}_{x}\backslash X$ et $X_{/x}$ sont les catégories
au-dessous et au-dessus de $x$.

\pagerange{1}{4}
\section*{1. L'outil : la catégorie $\widetilde{\pi}\backslash B$ (pages 1 à 4)}

Soit $\pi : B \to (\mathrm{Ens})$ un foncteur. En le composant avec le
foncteur $\mathfrak{P}^{*}$ qui à un ensemble associe l'ensemble ordonné de
ses parties non vides, on obtient
\[
  \widetilde{\pi} : B \longrightarrow (\mathrm{Ens})
  \longrightarrow (\mathrm{Ord}) \longrightarrow (\mathrm{Cat}),
\]
et l'on note $\widetilde{\pi}\backslash B$ la catégorie qui lui est associée
par la construction de Grothendieck : ses objets sont les couples $(b,E)$ avec
$b \in \mathrm{Ob}\,B$ et $E$ une partie non vide de $\pi(b)$, et
\[
  \mathrm{Hom}\bigl((b,E),(b',E')\bigr)
  = \{\, u : b \to b' \ \mid \ \widetilde{\pi}(u)(E) \subset E' \,\}.
\]
L'inclusion, et non l'égalité, est ce qui fait tout l'intérêt de la
construction : elle rend le foncteur d'oubli vers $B$ fidèle sans le rendre
plein, et c'est elle qui sera exploitée plus bas.\footnote{Le feuillet 12
insiste sur ce point en écrivant explicitement « et non
$\widetilde{\pi}(u)(E) = E'$ ».}

Posons $Z = \pi\backslash B$. Chaque fibre $Z_{b}$ s'identifie à $\pi(b)$, et
pour toute partie non vide $E$ de $\pi(b)$ on dispose dans $\widehat{Z}$ de
l'objet
\[
  \widetilde{E} = \coprod_{z \in E} z .
\]
On obtient ainsi le foncteur canonique
\[
  \widetilde{\pi}\backslash B \longrightarrow \widehat{Z},
  \qquad (b,E) \longmapsto \widetilde{E},
\]
qui est le personnage de ces quatre feuillets. Le calcul qu'il refait deux
fois, la seconde sous deux traits d'annulation, est celui du produit :
\[
  \widetilde{E} \times \widetilde{E}'
  = \coprod_{(z,z') \in E \times E'} z \times z' ,
\]
les produits étant pris dans $\widehat{Z}$, ce qui exhibe le produit de
$(b,E)$ et $(b',E')$ dans $\widetilde{\pi}\backslash B$ sous la forme
$\bigl(b \times b',\, p_{1}^{*}(E) \cap p_{2}^{*}(E')\bigr)$.

Ce que le feuillet 3 observe, et qui servira plus tard, est que
$\widetilde{\pi}\backslash B$ est \emph{cofibrée} sur $B$, de fibre en $b$
l'ensemble ordonné $\mathfrak{P}^{*}(\pi(b))$.

\pagerange{6}{15}
\section*{2. Le morphisme de topos $f$ et le lemme conjectural (pages 6 à 15)}

\pagerange{6}{10}
\subsection*{2.1 La construction}

Ici l'hypothèse du dossier est posée pour la première fois. Soit $X$ une
petite catégorie telle que
\[
  \textbf{(H)} \qquad
  \text{tout objet } 0\text{-connexe de } \widehat{X} \text{ est } 1\text{-connexe.}
\]
Soit $F$ un objet de $\widehat{X}$ et $U = \mathrm{Im}(F \to e)$ l'image du
morphisme vers l'objet final. Quitte à remplacer $X$ par $Y = X/U$, on peut
supposer $F \to e$ épimorphique, et $\widehat{Y} \simeq \widehat{X}/U$ est le
topos induit.

Soit $\widetilde{\Delta}$ la catégorie des ensembles finis non vides, réalisée
par les $\widetilde{\Delta}_{n} = [0,n] \cap \mathbf{Z}$. Le foncteur
\[
  \pi = \pi_{F} : \widetilde{\Delta}^{\,\circ} \longrightarrow (\mathrm{Ens}),
  \qquad I \longmapsto \pi_{0}(F^{I}),
\]
est un objet de $\widehat{\widetilde{\Delta}}$, et l'on pose
$Z = (\widetilde{\Delta}_{/\pi})^{\circ}$. Un objet de
$\widetilde{\Delta}_{/\pi}$ est un couple $(\widetilde{\Delta}_{n}, \alpha)$
où $\alpha$ est une composante connexe de $F^{\,n+1}$, et l'on définit
\[
  \varphi : (\widetilde{\Delta}_{/\pi})^{\circ} \longrightarrow \widehat{Y},
  \qquad (\widetilde{\Delta}_{n}, \alpha) \longmapsto \alpha \subset F^{\,n+1}.
\]
Sur les flèches : si $\rho : \widetilde{\Delta}_{n} \to
\widetilde{\Delta}_{m}$ vérifie $\pi(\rho)(\beta) = \alpha$, alors
$\rho^{*} : F^{\widetilde{\Delta}_{m}} \to F^{\widetilde{\Delta}_{n}}$
restreint à $\beta$ se factorise par $\alpha$, d'où
$\varphi(\rho) : \beta \to \alpha$.

Prolongé par continuité, $\varphi$ donne un foncteur
$\Phi : \widehat{Z} \to \widehat{Y}$ qui commute aux colimites quelconques. Il
est de plus exact, donc c'est l'image inverse d'un morphisme de topos
\[
  f : \widehat{Y} \longrightarrow \widehat{Z}.
\]

\pagerange{9}{10}
\subsection*{2.2 Asphéricité de $f$, et ce qu'elle donne}

Le morphisme $f$ est $W_{0}$-asphérique : pour tout objet $z$ de $Z$, l'image
inverse $f^{*}(z) = \Phi(z)$ est la composante connexe $\alpha$, donc
$0$-connexe. Il s'ensuit que $f^{*}(G)$ est $0$-connexe pour tout $G$
$0$-connexe de $\widehat{Z}$.\footnote{Le feuillet 9 justifie ce passage par
le fait que $f^{*}$ commute aux produits fibrés, « ou tout au moins aux
intersections de sous-objets » — la seconde forme suffit et c'est elle qui
est employée.}

Si l'on invoque maintenant l'hypothèse \textbf{(H)}, qui vaut aussi dans le
topos induit $\widehat{Y}$, on obtient que $f^{*}(G)$ est $1$-connexe dès que
$G$ est $0$-connexe. En considérant le morphisme induit
$\widehat{Y}/f^{*}(G) \to \widehat{Z}/G$, on conclut :

\begin{quote}
\textbf{Sous (H), tout objet $0$-connexe de $\widehat{Z}$ est $1$-connexe.}
\end{quote}

L'hypothèse s'est donc transportée de $X$ à $Z$, et $Z$ est une catégorie
d'origine simplicial. C'est là que le dossier voudrait conclure, et c'est là
qu'il s'arrête :

\begin{quote}
\textbf{Lemme (?)} — \emph{énoncé conjectural, non démontré.} Soient $\pi$ un
objet de $\widehat{\widetilde{\Delta}}$ et
$Z = (\widetilde{\Delta}_{/\pi})^{\circ}$. Si tout objet $0$-connexe de
$\widehat{Z}$ est $1$-connexe et si $\pi(\widetilde{\Delta}_{0})$ est
ponctuel, alors $\pi$ est l'objet final, c'est-à-dire que
$\pi(\widetilde{\Delta}_{n})$ est ponctuel pour tout $n$.
\end{quote}

Le point d'interrogation est de lui, et il n'y a pas de démonstration dans le
dossier.\footnote{Il écrit « Mais je conjecture que l'on devrait avoir
ceci ». Les deux corollaires qui suivent en dépendent entièrement, et sont
donc conditionnels ; la lecture les donne comme tels.} Sous cette réserve il
en tire deux conséquences. La première : si $F$ est $0$-connexe dans
$\widehat{X}$, ses puissances cartésiennes $F^{n}$ le sont aussi, en
particulier $F \times F$. La seconde : $X$ est localement cofiltrante.

C'est la première apparition de la thèse du dossier, et elle arrive comme
corollaire d'un lemme non démontré. Tout le reste du dossier peut se lire
comme la recherche d'une autre route vers le même énoncé.

\pagerange{11}{15}
\subsection*{2.3 La reprise en généralité}

Les feuillets 11 à 15 refont la construction pour une catégorie $A$
quelconque en place de $\widetilde{\Delta}$, avec
$Z = (A_{/\pi})^{\circ} \simeq \pi\backslash B$ où $B = A^{\circ}$, et
retrouvent la catégorie $\widetilde{\pi}\backslash B$ de la section 1 avec son
foncteur canonique vers $\widehat{Z}$. La question posée est celle des
conditions sur $A$ et $\pi$ qui rendent ce foncteur \emph{pleinement fidèle}
et compatible aux limites finies.

Un seul énoncé général est dégagé, et il est correct :

\begin{quote}
\textbf{Lemme 1.} Soit $C \to B$ une catégorie fibrée et $I$ une petite
catégorie telle que les limites de type $I$ existent dans $B$ et dans chaque
fibre $C_{b}$. Alors elles existent dans $C$, et $C \to B$ commute aux limites
de type $I$.
\end{quote}

Son corollaire, appliqué à $C = \widetilde{\pi}\backslash B$, demande que
cette catégorie soit non seulement cofibrée sur $B$ — ce qu'elle est
toujours — mais \emph{fibrée}, ce qui est une condition sur $\pi$.\footnote{La
condition est effectivement remplie : le feuillet 14 observe que
$u_{*} = \widetilde{\pi}(u)$ admet l'adjoint à droite
$E' \mapsto \pi(u)^{-1}(E')$, et que celui-ci préserve les parties non vides
dès que $\pi(u)$ est surjective. C'est cette dernière réserve, que le feuillet
porte entre parenthèses, qui est l'hypothèse véritable.} Le raisonnement
s'interrompt cinq lignes après le début du feuillet 15 et le reste de la
page est blanc.

\pagerange{17}{20}
\section*{3. Une présentation du groupe fondamental d'une catégorie (pages 17 à 20)}

\pagerange{17}{17}
\subsection*{3.1 Le groupe $\Delta(E)$}

Pour un ensemble $E$, soit $\Delta(E)$ le groupe engendré par des symboles
$[x,y]$, $(x,y) \in E^{2}$, soumis aux relations
\[
  [x,x] = 1, \qquad [x,y][y,x] = 1, \qquad [x,y][y,z] = [x,z].
\]
C'est le groupe libre sur $E$ privé de son point base : si $E$ est fini de
cardinal $n$ et $E = [1,n]$, les $g_{i} = [i,i+1]$ pour $1 \le i \le n-1$ en
sont une base libre, et l'on a $[i,j] = g_{i}g_{i+1} \cdots g_{j-1}$ pour
$i < j$. En particulier $\Delta(E)$ est libre de rang $n-1$.\footnote{Le
feuillet écrit « le groupe libre à \emph{(illisible)} générateurs » ; le rang
$n-1$ est ce que sa propre base $g_{1}, \ldots, g_{n-1}$ donne, et il le
confirme en énumérant les points $1, 2, \ldots, n$. Il note aussi la
réalisation $[x,y] = x^{-1}y$ dans le groupe libre $L(E)$, qui identifie
$\Delta(E)$ au noyau de $L(E) \to \mathbf{Z}$ envoyant chaque générateur sur
$1$.}

\pagerange{17}{17}
\subsection*{3.2 La présentation}

\begin{quote}
Soit $X$ une catégorie ayant un plus petit objet $e$ et un plus grand objet
$a$. Alors $\pi_{1}(X,e)$ est le quotient de $\Delta(\mathrm{Hom}(e,a))$ par
les relations
\[
  [\alpha v, \alpha u] = [\beta v, \beta u]
\]
pour tout système de doubles flèches $e \rightrightarrows x \rightrightarrows a$,
de flèches $u, v : e \to x$ et $\alpha, \beta : x \to a$.
\end{quote}

L'énoncé est celui du feuillet 17 et le dossier ne le démontre pas ; le
feuillet 18, biffé de deux diagonales, en porte les essais.

\pagerange{19}{20}
\subsection*{3.3 Le cas pseudo-cofiltrant}

C'est la forme dont la suite se servira, et elle est démontrée.

\begin{quote}
\textbf{Proposition 1.} Soit $X$ une catégorie pseudo-cofiltrante ayant un
plus petit objet $e$. Soient $M = \mathrm{End}_{X}(e)$, $X_{0}$ la
sous-catégorie pleine réduite à $e$, et $G = \pi_{1}(M) = \pi_{1}(X_{0}, e)$
le groupe enveloppant de $M$. Alors :

\textbf{a)} L'homomorphisme canonique $\pi_{1}(X_{0}, e) \to \pi_{1}(X,e)$ est
surjectif.

\textbf{b)} Pour toute double flèche $u, v : e \rightrightarrows x$ il existe
$u', v' \in M$ avec $uu' = vv'$, et l'élément
$[v,u] \overset{\text{déf}}{=} v'u'^{-1}$ de $G$ ne dépend pas du choix du
couple $(u',v')$ ; son image dans $\pi_{1}(X,e)$ est la classe du lacet
$e \xrightarrow{u} x \xleftarrow{v} e$.

\textbf{c)} Le noyau de la surjection $G \to \pi_{1}(X,e)$ est le sous-groupe
distingué engendré par les éléments
\[
  [w, gv]\,[v, fu]\,[w, gfu]^{-1}
\]
attachés aux diagrammes commutatifs
$x \xrightarrow{f} y \xrightarrow{g} z$ munis de $u : e \to x$, $v : e \to y$,
$w : e \to z$ compatibles.
\end{quote}

Le groupe fondamental d'une catégorie pseudo-cofiltrante à plus petit objet
est donc entièrement calculable à partir du monoïde $M$ et de ces relations.
C'est le théorème de structure sur lequel repose tout le massif suivant.

Il note enfin que si $e$ est aussi plus grand objet de $X$, toute relation du
type c) est triviale, de sorte que
$\pi_1(X_{0},e) \to \pi_{1}(X,e)$ est un isomorphisme.\footnote{Le feuillet
écrit l'indice $X_{0}$ des deux côtés de la dernière chaîne d'isomorphismes,
là où l'argument demande $\pi_{1}(X_{0},e) \xrightarrow{\sim} \pi_{1}(X,e)$.
La transcription le signale ; la forme correcte est celle donnée ici.}

\pagerange{21}{41}
\section*{4. Le contre-exemple : une catégorie $0$-connexe non $1$-connexe (pages 21 à 41)}

C'est le massif le plus long et le seul entièrement mené à terme. Le principe
en est simple : construire une catégorie dont le groupe fondamental se
surjecte sur $\mathbf{Z}$, en fabriquant à la main un « degré » à valeurs
entières.

\pagerange{20}{22}
\subsection*{4.1 Le dispositif}

Soient $M$ un monoïde, $\varphi, \psi \in M$ deux éléments fixés, et $R_{\varphi}$
la relation d'équivalence
\[
  u \sim_{R} v \iff \exists\, p, q \in \mathbf{N}, \quad \varphi^{p} u = \varphi^{q} v .
\]
Soient $Y = B_{M}$ le classifiant de $M$ — la catégorie à un objet dont le
monoïde des endomorphismes est $M$ —, $F = M/R_{\varphi}$ vu comme $M$-ensemble
à droite, donc comme objet de $\widehat{Y}$, et
\[
  X = Y_{/F}.
\]
Les objets de $X$ sont les éléments de $F$, et
$\mathrm{Hom}(x,y) = \{\, u \in M \mid x = yu \,\}$. L'objet $a = \dot{1}$ est
le plus grand.

Une première condition assure que $X$ a aussi un plus petit objet, de la forme
$e = a\psi$ :
\[
  \textbf{(A)} \qquad \forall\, u \in M,\ \exists\, u' \in M,
  \quad uu' \sim_{R} \psi ,
\]
condition remplie dès que $\psi$ est minimal dans $M$.

\pagerange{22}{24}
\subsection*{4.2 Le degré}

Posons $M_{0} = \mathrm{End}_{X}(e) = \{\, f \in M \mid \psi f \sim_{R} \psi \,\}$,
c'est-à-dire
\[
  f \in M_{0} \iff \exists\, p, q \in \mathbf{N}, \quad \varphi^{p}\psi f = \varphi^{q}\psi .
\]
\footnote{Le feuillet 22 trace ici le quantificateur universel, nettement
distinct du quantificateur existentiel qu'il forme au bas du feuillet 21. La
suite immédiate — « ne dépend pas du choix de $p,q$ » — n'aurait aucun sens
sous le quantificateur universel, qui rendrait d'ailleurs $M_{0}$ vide dès que
$\varphi$ n'est pas idempotent. C'est le quantificateur existentiel qui est
rétabli ici.}

Pour que le couple $(p,q)$ soit déterminé à translation près, on impose
\[
  \textbf{(B)} \qquad \varphi^{r}\psi = \varphi^{s}\psi \implies r = s .
\]
Alors $\delta(f) = q - p$ est bien défini et l'on obtient une application
\[
  \delta : M_{0} \longrightarrow \mathbf{Z}
\]
qui est un homomorphisme de monoïdes de $(M_{0}, \cdot)$ dans
$(\mathbf{Z}, +)$, c'est-à-dire $\delta(fg) = \delta(f) + \delta(g)$.\footnote{Le
feuillet 23 l'appelle « multiplicative » tout en écrivant la loi additive. Il
s'agit d'un homomorphisme vers le groupe additif $\mathbf{Z}$, et c'est ainsi
qu'il est nommé ici. La vérification qu'il donne est correcte et tient en deux
lignes : de $\varphi^{p}\psi f = \varphi^{q}\psi$ et
$\varphi^{p'}\psi g = \varphi^{q'}\psi$ on tire
$\varphi^{p+p'}\psi fg = \varphi^{q+q'}\psi$.} Comme $\mathbf{Z}$ est un
groupe, $\delta$ se factorise par le groupe enveloppant :
\[
  M_{0} \longrightarrow G_{0} = \pi_{1}(M_{0})
  \xrightarrow{\ \delta_{0}\ } \mathbf{Z}.
\]

La troisième condition demande que ce soit non trivial :
\[
  \textbf{(C)} \qquad \delta \neq 0 .
\]

\pagerange{25}{30}
\subsection*{4.3 Les relations s'annulent}

C'est le cœur du calcul, et il occupe six feuillets. Il s'agit de montrer
que $\delta_{0}$ tue les relations de la Proposition 1 c), donc descend en un
homomorphisme surjectif sur $\pi_{1}(X,e)$ lui-même.

Il commence par expliciter le crochet $[\mu,\lambda]$ pour une double flèche
$\lambda, \mu : e \rightrightarrows \xi = a\sigma$. Choisissant
$\lambda', \mu'$ avec $\lambda\lambda' = \mu\mu'$, puis $\nu$ tel que
$\sigma\lambda\lambda'\nu \sim_{R} \psi$, on obtient $\lambda'\nu, \mu'\nu \in M_{0}$
et
\[
  [\mu,\lambda] = (\mu'\nu)(\lambda'\nu)^{-1}.
\]
C'est ici qu'il renonce à la généralité, en toutes lettres : « Je crois que je
vais craquer, et restreindre la généralité ». Il suppose désormais $\psi$
\emph{minimal}, de sorte qu'il existe $\theta(u)$ avec $u\,\theta(u) = \psi$ —
égalité stricte et non simple congruence — et prend
$\lambda' = \theta(\lambda)$, $\mu' = \theta(\mu)$.

Appliqué aux trois crochets de la relation fondamentale, ce calcul donne
quatre éléments de $M_{0}$ :
\[
  c_{1} = \theta(gv)\,\theta(\rho\psi), \quad
  c_{2} = \theta(v)\,\theta(\rho g\psi), \quad
  c_{3} = \theta(fu)\,\theta(\rho g\psi), \quad
  c_{4} = \theta(gfu)\,\theta(\rho\psi),
\]
et la relation à annuler est $c_{1}c_{2}^{-1}c_{3}c_{4}^{-1} = 1$.\footnote{Le
feuillet 28 porte, sur le dernier facteur, un exposant $-1$ placé juste après
la parenthèse fermant $\theta(gfu)$, et la parenthèse ouvrante du facteur
n'est jamais refermée. Sa propre règle de composition, posée deux feuillets
plus haut, donne $\bigl(\theta(gfu)\,\theta(\rho\psi)\bigr)$ sans exposant, et
c'est cette forme qui est retenue. Il remarque lui-même, avec un certain
étonnement, que « $w$ y a disparu sans laisser de traces » : c'est correct, et
la raison en est qu'on peut poser $w = gv$ sans restreindre la généralité.}

Une quatrième condition intervient alors :
\[
  \textbf{(D)} \qquad \psi \sim_{R} 1 ,
\]
qui permet de récrire les contraintes sous la forme
$\rho g f u, \rho g v, \rho w \in M_{0}$. Posant
$\pi = \rho g f u$ et $\pi' = \rho g v$, on obtient
\[
  \pi c_{1} = \pi c_{2} = \psi, \qquad \pi' c_{3} = \pi' c_{4} = \psi,
\]
d'où $\delta(c_{1}) = \delta(c_{2})$ et $\delta(c_{3}) = \delta(c_{4})$, donc
\[
  \delta_{0}\bigl(c_{1}c_{2}^{-1}c_{3}c_{4}^{-1}\bigr) = 0 .
\]
Le degré tue donc toutes les relations, et se factorise en
\[
  \delta_{0}' : \pi_{1}(X,e) \twoheadrightarrow \mathbf{Z}.
\]

\pagerange{31}{33}
\subsection*{4.4 L'énoncé}

\begin{quote}
\textbf{Proposition 2.} Soient $M$ un monoïde et $\varphi, \psi \in M$ tels
que :

\textbf{a)} $\psi$ est minimal dans $M$, c'est-à-dire que pour tout $u \in M$
il existe $\theta(u)$ avec $\psi = u\,\theta(u)$ ;
\textbf{b)} $\psi \sim_{R_{\varphi}} 1$ ;
\textbf{c)} $\varphi^{r}\psi = \varphi^{s}\psi$ entraîne $r = s$.

Soient $Y = B_{M}$, $F = M/R_{\varphi}$ et $X = Y_{/F}$. Alors $e = \dot{1}$
est à la fois plus petit et plus grand objet de $X$, et il existe un
homomorphisme surjectif $\pi_{1}(X,e) \to \mathbf{Z}$ ; en particulier $X$
n'est pas $1$-connexe.
\end{quote}

L'hypothèse a) entraîne que $M$, donc $Y = B_{M}$, est pseudo-cofiltrant, ce
qu'il note en marge de l'énoncé.

\begin{quote}
\textbf{Corollaire.} Si $M$ est isomorphe à $\mathrm{Ep}(E)$ ou à
$\mathrm{Mon}(E)$ pour $E$ infini, alors $M$ est pseudo-cofiltrant,
$Y = B_{M}$ est $1$-connexe, et il existe $F$ dans $\widehat{Y}$ tel que
$X = Y_{/F}$ soit $0$-connexe et non $1$-connexe.
\end{quote}

Pour l'appliquer il suffit de prendre $\psi = \varphi$ un élément minimal
quelconque, et de vérifier la condition c), ce qu'il fait par l'absurde : si
$\varphi^{r} = \varphi^{s}$ avec $r = s+h$, $h > 0$, alors
$\varphi^{h}\varphi^{s} = \varphi^{s}$, donc $\varphi$ serait inversible. Or un
élément inversible ne peut être minimal, sans quoi tout élément de $M$ le
serait et $M$ serait un groupe.\footnote{L'argument, donné au feuillet 33, est
joli et mérite d'être suivi : de $\varphi = uu'$ et $v = u'\varphi^{-1}$ on
tire $uv = 1$, puis $p = vu$ vérifie $p^{2} = p$, d'où $p = 1$ par
simplification à droite. Il est donné là dans le cas particulier des monoïdes
où la simplification est licite ; le feuillet 35 le reprend en toute
généralité dans le quotient $\overline{M}$, où elle l'est par construction.}

C'est aussi à ce feuillet qu'il note la difficulté qu'il a eue :
« j'ai un peu mal à construire un $F$ $0$-connexe qui ne soit $1$-connexe ! ».

\pagerange{35}{37}
\subsection*{4.5 Le quotient simplifiable}

Les feuillets 35 à 37 dégagent l'outil qui rend l'argument précédent robuste.
Si $\psi$ est minimal, la relation $u \sim_{R_{0}} v \iff u\psi = v\psi$ est
celle définie par l'application
\[
  M \longrightarrow \mathrm{Ep}(M_{0}), \qquad
  u \longmapsto (\lambda \mapsto u\lambda),
\]
et le quotient $\overline{M} = M/R_{0} \subset \mathrm{Ep}(M_{0})$ est un
monoïde \emph{simplifiable à droite} — $\bar{u}\bar{f} = \bar{v}\bar{f}$
entraîne $\bar{u} = \bar{v}$ — et pseudo-cofiltrant comme image de $M$. La
condition c) de la Proposition 2 se lit alors
$\overline{\varphi}^{\,r} = \overline{\varphi}^{\,s} \implies r = s$.

De là l'alternative qui organise la fin du massif. Ou bien
$\overline{\varphi}$ est d'ordre fini, donc inversible, et l'on montre que
$\overline{M}$ est un groupe ; ou bien la condition c) est satisfaite et la
Proposition 2 s'applique.

\begin{quote}
\textbf{Proposition 3} (énoncé incomplet). Soit $M$ un monoïde ayant un
élément minimal $\psi$, et $\varphi \in M$ tel que $\varphi^{\beta}$ soit
minimal pour un $\beta$. S'il existe $r \neq s$ avec
$\varphi^{r}\psi = \varphi^{s}\psi$, l'énoncé s'interrompt ici.
\end{quote}

L'énoncé s'arrête là ; le bas du feuillet porte une grande boucle et la
conclusion manque. Ce que le feuillet 37 en tire, en revanche, est complet :
dans le cas où $\overline{M}$ est un groupe, si $Y = B_{M}$ est $1$-connexe ce
groupe est trivial, donc $u\lambda = \lambda$ pour tout $u$ et tout $\lambda$
minimal, et $M$ est cofiltrant. Dans le cas contraire, $Y_{/F}$ est
$0$-connexe, non $1$-connexe, de groupe fondamental $\pi_{1}(M_{0})$, lequel
se surjecte sur $\mathbf{Z}$ ; donc $H^{1}(X,\mathbf{Z}) \neq 0$.\footnote{Il
vérifie au passage que $\delta$ est bien non nul : comme $\psi \in M_{0}$, on a
$\varphi\psi \in M_{0}$, donc il existe $f$ avec $\psi f = \varphi\psi$, et un
tel $f$ vérifie $\delta(f) = 1$.}

\pagerange{39}{41}
\subsection*{4.6 Le théorème, et son interruption}

\begin{quote}
\textbf{Corollaire.} Soit $M$ un monoïde ayant un élément minimal. Si tous les
objets $0$-connexes de $\widehat{B_{M}}$ sont $1$-connexes — ou même seulement
si leur $H^{1}(-,\mathbf{Q})$ est nul — alors $M$ est cofiltrant.
\end{quote}

C'est enfin la thèse du dossier, obtenue cette fois sans lemme conjectural,
mais seulement pour les monoïdes. Le feuillet 39 entreprend de la transporter
aux catégories :

\begin{quote}
\textbf{Théorème} (énoncé interrompu). Soit $X$ une catégorie
pseudo-cofiltrante dont les objets $0$-connexes de $\widehat{X}$ vérifient
l'hypothèse de connexité, admettant un plus petit objet $e$ tel que
$M = \mathrm{End}_{X}(e)$ ait un élément minimal $\varphi$, l'énoncé s'interrompt ici.
\end{quote}

La démonstration se sépare en deux cas selon que $\mathbf{N}$ opère sur $e$
par $\varphi$ sans point fixe ou non, et s'interrompt à l'intérieur du premier.

Le feuillet 41 reprend exactement là.\footnote{C'est le feuillet dont la
pagination ne se raccorde pas : il porte « $7$ » de sa main alors que la suite
précédente est numérotée $3$ à $18$. Le raccord du sens, lui, est net, et
c'est celui qui est suivi ici.} Dans le premier cas,
$\varphi^{p} = \varphi^{q}$ entraîne $p = q$, et l'objet quotient
$F = e/\mathbf{N}$ est $0$-connexe de $\pi_{1}$ se surjectant sur $\mathbf{Z}$,
donc non $1$-connexe. Dans le second, il existe $p \neq q$ avec
$\varphi^{p} = \varphi^{q}$ ; l'image de $M$ dans $\mathrm{Ep}(E_{0})$, où
$E_{0}$ est l'ensemble des éléments minimaux, est alors un groupe $\Gamma$, et
\[
  \pi_{1}(X_{0}, e) \simeq \pi_{1}(M) \twoheadrightarrow \Gamma .
\]
\footnote{Le feuillet écrit « l'ensemble $\mathrm{Ep}(E)$ des épimorphismes de
$E$ $=$ ensemble des éléments minimaux de $E$ », ce qui définit $E$ par
lui-même. L'argument demande deux ensembles distincts, et le second reçoit le
nom $E_{0}$ quatre feuillets plus loin ; c'est cette lecture qui est adoptée
ici.} Si $X_{0}$ est $1$-connexe, alors $\Gamma$ est trivial, donc
$u\varphi = \varphi$ pour tout $u$, et $X_{0}$ est cofiltrant. D'où sa
conclusion, écrite en toutes lettres :

\begin{quote}
\textbf{Conclusion.} Si tout objet de $\widehat{X}$ est $1$-connexe, alors $X$
est cofiltrant.
\end{quote}

\pagerange{44}{53}
\section*{5. Le cas fini, et sa correction (pages 44 à 53)}

\pagerange{44}{46}
\subsection*{5.1 Les monoïdes finis}

\begin{quote}
\textbf{Proposition 1.} Soit $M$ un monoïde \emph{fini} dont le groupe
enveloppant est trivial et qui est pseudo-cofiltrant. Alors il existe
$p \in M$ tel que $up = p$ pour tout $u \in M$ — en particulier
$p^{2} = p$ — et $M$ est cofiltrant.
\end{quote}

La démonstration est complète et c'est le plus solide résultat du dossier.
Elle tient en trois pas. D'abord, la pseudo-cofiltrance donne pour toute
famille finie $(u_{i})$ un $p$ qui se factorise par chaque $u_{i}$ ; $M$ étant
fini, il existe un $p$ tel que $p = u\,v(u)$ pour \emph{tout} $u \in M$.
Ensuite, on réalise $M$ comme monoïde d'endomorphismes d'un ensemble fini $E$,
par exemple par la représentation régulière. La factorisation donne
$\mathrm{Im}\,p \subset \mathrm{Im}\,u$ pour tout $u$, donc
$E_{0} = \mathrm{Im}\,p$ est le plus petit des $\mathrm{Im}\,u$ ; et pour tout
$u$,
\[
  u(E_{0}) = u\,p(E) = \mathrm{Im}(up) \supset E_{0},
\]
d'où $u(E_{0}) = E_{0}$ puisque $E_{0}$ est fini.\footnote{La finitude de
$E_{0}$ est ce qui transforme l'inclusion en égalité, et le feuillet l'emploie
sans la rappeler à cet endroit ; elle est bien dans les hypothèses. Sans elle
l'argument tombe, et c'est exactement ce qui se passe pour
$\mathrm{Ep}(E)$ avec $E$ infini, étudié à la section 6.} La restriction
$u \mapsto u|E_{0}$ est donc un homomorphisme $M \to \mathrm{Aut}(E_{0})$ vers
un groupe ; l'hypothèse sur le groupe enveloppant le rend trivial, d'où
$u\,p(x) = p(x)$ pour tout $x$, c'est-à-dire $up = p$.

En langage moderne, $p$ est un idempotent qui engendre un idéal à gauche
minimal réduit à un zéro à gauche : c'est le noyau de Rees du monoïde fini,
dont la trivialité du groupe enveloppant force qu'il soit ponctuel.\footnote{La
théorie de structure des semi-groupes finis — noyau, semi-groupe complètement
simple — était disponible ; rien sur ces feuillets n'indique qu'il l'ait eue
en vue, et la démonstration est menée de bout en bout par ses propres moyens.}

\pagerange{46}{49}
\subsection*{5.2 La transposition aux catégories, et son échec}

Il propose ensuite, avec deux points d'interrogation de sa main :

\begin{quote}
\textbf{Proposition 2 (??)} Soit $X$ une catégorie finie, $1$-connexe,
pseudo-cofiltrante et ayant un plus grand objet. Alors $X$ est cofiltrante.
\end{quote}

Les feuillets 47 et 48 tentent de calquer la démonstration précédente en
choisissant, pour chaque objet $b$, une flèche $p_{b} : e \to b$ compatible
avec un $p_{a}$ fixé. Le point d'achoppement est l'unicité : $p_{b}$ n'est pas
déterminé, et le feuillet 48 se réduit à des diagrammes repris deux fois sans
prose suivie.

Le feuillet 49 en tire néanmoins un corollaire, soigneusement rédigé, qui
déduit le cas des catégories finies du cas des monoïdes en passant par
l'ouvert $U$ des plus petits objets et l'équivalence $U_{0} \simeq B_{M}$.
Puis, en marge, de sa main :

\begin{quote}
\emph{« Faux ! revoir la démonstration »}
\end{quote}

Le corollaire est donc répudié par son auteur et n'est pas reproduit ici comme
un résultat.\footnote{La lecture ne cherche pas à réparer la démonstration ni
à décider ce qui y manque : la mention marginale est un fait du dossier, et
l'énoncé rétabli aux feuillets 50 à 53 est la réponse qu'il lui a lui-même
donnée. On notera seulement que le corollaire repose sur le passage de $X$ à
l'ouvert cocofinal $U$ puis à $U_{0}$, et que c'est la préservation de la
$1$-connexité le long de ce passage qui est le point douteux.} Il note
cependant, et cela reste valable, qu'au lieu de la finitude de $X$ il suffit
de supposer que $X$ a un plus petit objet $e$ et que $\mathrm{End}_{X}(e)$ est
fini.

\pagerange{50}{53}
\subsection*{5.3 La seconde mouture : les systèmes locaux $H_{x}$}

Les feuillets 50 à 53 reprennent la proposition depuis le début, sur d'autres
bases. Ce n'est pas une répétition : l'objet porteur change, et c'est le
changement qui est l'information.

Soit $X$ pseudo-cofiltrante à plus petit objet $e$, $M = \mathrm{Hom}(e,e)$,
et $F(x) = \mathrm{Hom}(e,x)$. On suppose que les sous-$M$-ensembles de $F(x)$
satisfont la condition minimale — ce qui vaut si les $F(x)$ sont finis. Alors
parmi les sous-$M$-ensembles de la forme $fM$ il en existe un plus petit, et
c'est
\[
  H_{x} = \bigcap_{f \in F(x)} fM .
\]
Les $H_{x}$ forment un système local sur $X$ : pour $f : x \to y$ on a
$H_{y} = f \cdot H_{x}$, et si $e$ est aussi objet maximal, l'application
$\lambda \mapsto f\lambda$ est une bijection $H_{x} \xrightarrow{\sim} H_{y}$.

La $1$-connexité de $X$ entre alors en jeu exactement là où il faut : elle
force ce système local à être \emph{constant}, donc l'application
$u \mapsto fu$ de $H_{e}$ dans $H_{x}$ ne dépend pas du choix de
$f \in F(x)$, c'est-à-dire
\[
  fu = gu \qquad \text{pour } u \in H_{e},\ f, g : e \rightrightarrows x .
\]
C'est précisément l'égalisation qui manque à la pseudo-cofiltrance pour faire
une cofiltrance. D'où :

\begin{quote}
\textbf{Proposition 2} (seconde mouture). Soit $X$ une catégorie finie
pseudo-cofiltrante et $X_{0}$ l'ouvert de ses plus petits objets, qui est non
vide et $0$-connexe. Si $X_{0}$ est $1$-connexe, alors $X$ est cofiltrante —
donc tout objet $0$-connexe de $\widehat{X}$ est $\infty$-connexe, et plus
généralement $\infty$-asphérique pour tout localisateur fondamental $W$
convenable.
\end{quote}

\begin{quote}
\textbf{Corollaire.} Soit $X$ une catégorie finie telle que tout objet
$0$-connexe de $\widehat{X}$ soit $1$-connexe — par exemple telle que $X_{0}$
soit totalement $W$-asphérique pour un localisateur fondamental
$W \subset W_{1}$. Alors $X$ est cofiltrante, $X^{\circ}$ est filtrante, et
tout objet de $\widehat{X}$ est $\infty$-connexe.
\end{quote}

C'est le premier endroit du dossier où le vocabulaire des \emph{localisateurs
fondamentaux} apparaît en toutes lettres, et il y arrive comme conséquence, non
comme cadre.\footnote{L'ouvert des plus petits objets est noté $X_{0}$ dans la
proposition et $X^{0}$, l'indice porté en exposant, dans le corollaire ; une
seule notation est employée ici. La condition « $W$ satisfaisant $W(\delta)$ »
que porte la proposition n'est pas lisible avec certitude sur le feuillet et
n'est pas reconstruite.}

\pagerange{55}{58}
\section*{6. Deux monoïdes infinis (pages 55 à 58)}

Ces feuillets fournissent la matière première du contre-exemple : des monoïdes
pseudo-cofiltrants, à groupe enveloppant trivial, et \emph{non} cofiltrants.
Ils montrent du même coup que la finitude est indispensable à la Proposition 1
de la section 5.

Soit $E$ un ensemble infini et $M = \mathrm{Ep}(E)$ le monoïde de ses
endomorphismes surjectifs. Alors :

\begin{enumerate}
\item[$1^{\circ}$] $uf = vf \implies u = v$ : le monoïde est simplifiable à
droite, donc \emph{non} cofiltrant — une égalisation forcerait $u = v$.
\item[$2^{\circ}$] $M$ est pseudo-cofiltrant et possède des éléments
minimaux : ce sont les surjections dont toutes les fibres ont le cardinal de
$E$. Mais on n'a jamais $uw = w$, ni $u^{2} = u$ pour $u \neq 1$.
\item[$3^{\circ}$] Les éléments minimaux sont exactement les $w_{0}\sigma$
avec $\sigma \in M^{*} = \mathrm{Aut}(E)$.
\item[$4^{\circ}$] Il existe $w$ minimal et $f$ avec $fw = f$, donc l'image
$\overline{w}$ de $w$ dans le groupe enveloppant $G$ est $1$.
\item[$5^{\circ}$] L'ensemble des $\sigma \in M^{*}$ tels que $f\sigma = f$
pour un certain $f$ engendre le groupe $M^{*}$.
\item[$6^{\circ}$] Par conséquent le groupe enveloppant $G$ de $M$ est
trivial.
\end{enumerate}

L'enchaînement de $5^{\circ}$ et $6^{\circ}$ est celui-ci : l'image de $M^{*}$
dans $G$ est triviale par $5^{\circ}$ et $4^{\circ}$, donc par $3^{\circ}$
tout élément minimal a une image triviale, et par $2^{\circ}$ tout $u$
s'écrit avec des minimaux de part et d'autre, d'où $\overline{u} = 1$.

Le monoïde $\mathrm{Mon}(E)$ des endomorphismes injectifs satisfait les
propriétés duales.\footnote{Le passage de l'un à l'autre est le passage au
monoïde opposé, qu'il écrit $M^{\circ}$. La propriété duale de
$1^{\circ}$ est la simplification à gauche.}

Le feuillet 57 ajoute un lemme qui servira à la section suivante, et qui est
la raison pour laquelle la notation $X_{/\!/a}$ est celle du cocrible :

\begin{quote}
\textbf{Lemme.} Soit $X$ localement irréductible. Alors la relation
$R(a,b) : X_{/\!/a} \cap X_{/\!/b} \neq \emptyset$ est une relation
d'équivalence sur $\mathrm{Ob}\,X$, et c'est celle définie par
$\mathrm{Ob}\,X \to \pi_{0}(X)$. De plus, si $X_{/\!/a} \cap X_{/\!/b}$ est
non vide, il est $0$-connexe. Pour que $X$ soit $1$-connexe, il suffit que
les $X_{/\!/a}$ le soient.
\end{quote}

\pagerange{60}{79}
\section*{7. Localisateurs, asphéricité et suspensions (pages 60 à 79)}

La dernière campagne du dossier reprend la question par la géométrie. Elle est
aussi celle où le vocabulaire des dérivateurs est le plus présent.

\pagerange{70}{70}
\subsection*{7.1 Le critère d'asphéricité totale}

\begin{quote}
\textbf{Proposition.} Pour que $X^{\circ}$ soit totalement $W$-asphérique il
suffit que tout objet $0$-connexe de $\widehat{X}$ soit $W$-asphérique ; et
c'est aussi nécessaire si $W$ satisfait $\mathrm{Loc}(6)$.
\end{quote}

La démonstration est en deux temps. Dans un sens, $X_{/e} \simeq X_{/F}$ et
l'asphéricité de $F$ équivaut à la cocofinalité de l'inclusion
$X_{/F} \hookrightarrow X'$. Dans l'autre, si $X \to Y$ est cocofinal alors
les $X_{/y}$ sont $0$-connexes, donc $W$-asphériques par hypothèse, donc
$X \to Y$ est $W$-asphérique et $Y$ est $W$-asphérique.

\begin{quote}
\textbf{Corollaire.} Si $W$ satisfait $\mathrm{Loc}(8)$ et si $X$ est
cofiltrante, les objets $0$-connexes de $\widehat{X}$ sont $W$-asphériques,
donc $X$ est totalement $W$-asphérique.
\end{quote}

\footnote{Les axiomes $\mathrm{Loc}(6)$, $\mathrm{Loc}(6a)$ et
$\mathrm{Loc}(8)$ sont invoqués par leurs numéros et ne sont énoncés nulle part
dans le dossier. Cette lecture ne leur invente pas d'énoncé et les cite comme
il les cite. Le lecteur qui voudra les retrouver les cherchera du côté de la
liste d'axiomes des localisateurs fondamentaux.}

\pagerange{72}{74}
\subsection*{7.2 Le dispositif géométrique}

Soit $S$ la catégorie ordonnée à quatre objets
\[
  S = \bigl(\, a \to \xi \to b \ ; \ a \to \xi' \to b \,\bigr)
\]
— deux chemins parallèles de $a$ à $b$ — et soit $Z$ obtenue en lui adjoignant
un objet final $e$. Un foncteur $S \to X$ étant donné, on forme $X'$ par le
carré cocartésien

\begin{tikzcd}
S \arrow[r, hook] \arrow[d] & Z \arrow[d] \\
X \arrow[r, hook] & X'
\end{tikzcd}

Le calcul central est celui des flèches nouvelles. Pour $x$ dans $X$,
\[
  \mathrm{Hom}_{X'}(x,e) \simeq \pi_{0}\bigl({}_{x}\backslash S\bigr),
\]
cette colimite étant celle du losange

\begin{tikzcd}[column sep=small, row sep=small, nodes={font=\scriptsize}]
& \mathrm{Hom}(x,\xi) & \\
\mathrm{Hom}(x,a) \arrow[ur] \arrow[dr] & & \mathrm{Hom}(x,b) \arrow[ul] \arrow[dl] \\
& \mathrm{Hom}(x,\xi') &
\end{tikzcd}
c'est-à-dire
\[
  \mathrm{Hom}_{X'}(x,e) \simeq E(x)/R(x),
  \qquad E(x) = \mathrm{Hom}(x,\xi) \amalg \mathrm{Hom}(x,\xi'),
\]
où $R(x)$ est la relation d'équivalence engendrée par
\[
  (*) \qquad
  \begin{cases}
    \lambda v \sim \lambda' v & \text{pour } v : x \to a, \\
    \mu w \sim \mu' w & \text{pour } w : x \to b,
  \end{cases}
\]
$\lambda, \lambda'$ désignant les deux flèches issues de $a$ et $\mu, \mu'$
les deux flèches aboutissant de $b$.\footnote{Le feuillet 77 mène le même
calcul avec la catégorie plus simple $D = (a \rightrightarrows b)$ et conclut
que $X \to X'$ est cofinal. C'est une mouture antérieure, gardée distincte
ici : le passage de $D$ à $S$ est ce qui fait apparaître les deux composantes
$\xi$ et $\xi'$, donc la somme $E(x)$, donc tout le reste.}

\pagerange{76}{78}
\subsection*{7.3 Le lemme de recouvrement}

Posons $\tilde X = X_{/e}$ et notons $\tilde a$, $\tilde b$, $\tilde\xi$,
$\tilde\xi'$ les relevés.\footnote{Le feuillet 66 \emph{définit} au contraire
$\tilde X = e \backslash X$. Les deux feuillets ne s'accordent pas ; la
transcription porte chaque fois ce qui est tracé, et cette lecture retient
$X_{/e}$, qui est la définition du feuillet 76 où le lemme est énoncé et
démontré.}

\begin{quote}
\textbf{Lemme.} On a
\[
  \tilde X = \tilde X_{/\!/\tilde\xi} \ \cup \ \tilde X_{/\!/\tilde\xi'},
  \qquad
  \tilde X_{/\!/\tilde\xi} \cap \tilde X_{/\!/\tilde\xi'}
  = \tilde X_{/\!/\tilde a} \ \cup \ \tilde X_{/\!/\tilde b}.
\]
\end{quote}

La première égalité est immédiate. La seconde est démontrée au feuillet 78, et
la démonstration est exactement le maniement de la relation $R(x)$. Soient
$u \in \mathrm{Hom}(x,\xi)$ et $u' \in \mathrm{Hom}(x,\xi')$ congrus modulo
$R(x)$ avec $u \neq u'$. Il existe alors une chaîne
$f_{0} = u, f_{1}, \ldots, f_{n} = u'$ dans $E(x)$, avec $n \ge 1$, dont
chaque maillon est l'une des deux relations $(*)$. En particulier, pour
$i = 0$, ou bien il existe $v : x \to a$ avec $f_{0} = \lambda v$, ou bien il
existe $w : x \to b$ avec $f_{0} = \mu w$. Donc $u$ se factorise par $a$ ou
par $b$, c'est-à-dire que $\tilde x$ appartient à
$\tilde X_{/\!/\tilde a} \cup \tilde X_{/\!/\tilde b}$.\footnote{La dernière
phrase du feuillet écrit $\tilde\xi$ là où le sujet doit être $\tilde x$ ;
c'est un lapsus manifeste et la forme correcte est donnée ici.}

\pagerange{61}{79}
\subsection*{7.4 Le calcul de suspensions}

Tout est alors en place. Posons
\[
  T = \tilde X_{/\!/(\tilde a, \tilde b)}
  = \tilde X_{/\!/\tilde a} \cap \tilde X_{/\!/\tilde b},
\]
la partie de $\tilde X$ que $\tilde a$ et $\tilde b$ atteignent tous deux. Le
lemme précédent dit que $\tilde X$ est recouvert par deux morceaux dont
l'intersection se décrit par $T$, et c'est la situation d'un recollement de
type Mayer-Vietoris. On obtient la tour
\[
  \tilde X \simeq \Sigma^{1}_{W}(T), \qquad
  X' \simeq \Sigma^{1}_{W}(X) \simeq \Sigma^{2}_{W}(\tilde X)
  \simeq \Sigma^{3}(T),
\]
où $\Sigma$ est la suspension.\footnote{L'argument du terme médian est tracé
d'un seul jambage sur le feuillet 79 et la transcription le donne comme
douteux ; c'est la chaîne $\tilde X \simeq \Sigma^{1}(T)$ et
$X \simeq \Sigma^{1}(\tilde X)$ du feuillet 66 qui rend seule cohérente la
lecture retenue.} Le feuillet 62 avait déjà noté le premier étage sous la
forme $Z = e\backslash X \sim S^{1}$ et $X' \sim S^{2}$.

Le corollaire du feuillet 79 en tire la conclusion vers laquelle tout le
dossier tendait :

\begin{quote}
\textbf{Si $T = \emptyset$, alors $X' \sim_{W_{\infty}} \Sigma^{2}(E) = S^{2}$.}
Par conséquent, si $W$ est plus fin que $H^{2}(-,k)$ pour un groupe abélien
$k$, alors $X'$ n'est pas $W$-asphérique. Donc, sous les hypothèses du
dossier, on doit avoir $T \neq \emptyset$.
\end{quote}

Ici $E = \{\beta,\gamma\}$ est la catégorie discrète à deux objets, qui est la
sphère $S^{0}$ ; sa double suspension est $S^{2}$. L'énoncé dit donc : le
dispositif fabrique, à partir d'un trou, une $2$-sphère, et une $2$-sphère est
un témoin qui distingue les localisateurs. Le même phénomène avait été noté
au feuillet 61 sous une forme purement catégorique :

\begin{quote}
S'il existe $B, C$ dans $\widehat{X}$ avec $B \times C = \emptyset$ et
$\beta, \gamma$ dans $X^{\vee}$ avec $\beta \times \gamma = \emptyset$, alors
$X$ n'est pas $1$-connexe.
\end{quote}

\footnote{La variance est ici décisive et elle est portée par le seul signe
$\wedge$ ou $\vee$ : le premier produit est celui de deux préfaisceaux, le
second celui de deux foncteurs covariants représentables. Prendre les deux du
même côté rendrait l'énoncé faux. Les huit dernières lignes du feuillet 61
sont, de l'aveu de la transcription, le passage le plus incertain du lot.}

Et le feuillet 68 en donne la forme quantitative : si $X$ est $0$-connexe sans
être filtrante, il existe $X \to X'$ cofinal avec
$H^{2}(X') \simeq H^{2}(S^{2})$, de sorte qu'aucun localisateur $W \subset W_{2}$
ne peut rendre $X$ totalement asphérique.

\pagerange{81}{83}
\section*{8. Là où le dossier s'arrête : congruence contre égalité (pages 81 et 83)}

Les deux derniers feuillets écrits cherchent à décider si $T$ peut être vide,
c'est-à-dire à interpréter $\tilde X_{/\!/(\tilde a,\tilde b)}$ directement.
Ils échouent, et l'intérêt est dans la façon dont ils échouent.

Soient $v : x \to a$ et $w : x \to b$. Les quatre composés
$\lambda v$, $\mu w$, $\lambda' v$, $\mu' w$ sont des éléments de $E(x)$, et la
question est de savoir s'ils sont tous congrus modulo $R(x)$. Deux des trois
congruences sont acquises par la définition même de $R(x)$ :
\[
  \lambda v \sim \lambda' v, \qquad \mu w \sim \mu' w .
\]
Ce qui manque est le pont entre les deux familles :
\[
  \boxed{\ \lambda v \sim \mu w\ }
\]
et c'est sur cette formule encadrée que s'arrête le feuillet 81.\footnote{Le
feuillet écrit les deux sources sous les accolades différemment — un $m$ d'un
côté, un $a$ de l'autre — alors que la marche du raisonnement demande la même
source de part et d'autre. Le feuillet 74 écrit la même somme avec une source
unique, $\mathrm{Hom}(x,\xi) \amalg \mathrm{Hom}(x,\xi')$, et c'est cette
forme qui est employée ici ; la transcription, elle, garde le tracé.}

Le feuillet 83 donne la réponse, et elle est négative :

\begin{quote}
De $\lambda' v = \mu' w$ on conclut $\lambda v \sim \mu w$,
\emph{« mais nullement $\lambda v = \mu w$ »}.
\end{quote}

Toute la difficulté du dossier tient dans cette phrase. La congruence modulo
$R(x)$ suffit à peupler $\tilde X_{/\!/(\tilde a,\tilde b)}$ — elle donne
$T \neq \emptyset$ — mais elle ne suffit pas à produire l'objet $\tilde x$ qui
réaliserait la convergence au sens strict. L'écart entre $\sim$ et $=$ est
exactement l'écart entre pseudo-cofiltrance et cofiltrance, ramené ici à sa
plus simple expression, et le dossier le reconnaît sans le franchir. Ses
derniers mots :

\begin{quote}
\emph{« mais je ne vois pas de moyen de montrer en même temps que l'autre le
$\tilde x$. »}
\end{quote}

\footnote{Il écrit le signe $=$ pour le mot \emph{même} — « en $=$ temps » —
abréviation qu'il emploie ailleurs dans le fonds. Le feuillet, et avec lui le
dossier, s'arrête sur cette phrase ; les feuillets 84 et 85 sont les deux
faces d'une enveloppe.}

\pagerange{1}{85}
\section*{Ce que le dossier établit, et ce qu'il laisse}

Il vaut la peine de séparer les deux, parce qu'un dossier de travail ne
distingue pas et qu'une lecture doit le faire.

\textbf{Établi, avec démonstration complète :} la Proposition 1 du feuillet 44
sur les monoïdes finis ; la Proposition 1 des feuillets 19 et 20 sur la
présentation de $\pi_{1}$ d'une catégorie pseudo-cofiltrante ; la
Proposition 2 du feuillet 31 et son corollaire, qui fournissent le
contre-exemple ; les six propriétés de $\mathrm{Ep}(E)$ du feuillet 55 ; le
lemme de recouvrement du feuillet 76, démontré au feuillet 78 ; la proposition
d'asphéricité totale du feuillet 70 ; la seconde mouture des feuillets 50
à 53.

\textbf{Conjectural ou interrompu :} le Lemme (?) du feuillet 10, dont
dépendent ses deux corollaires ; la Proposition 2 (??) du feuillet 46,
abandonnée ; la Proposition 3 du feuillet 35, dont l'énoncé s'arrête avant sa
conclusion ; le Théorème du feuillet 39, interrompu dans son premier cas ; le
corollaire du feuillet 49, répudié par lui en marge ; et la question finale
des feuillets 81 et 83, laissée ouverte.

\textbf{Ce que le dossier ne fait pas :} il n'énonce nulle part les axiomes
$\mathrm{Loc}(6)$, $\mathrm{Loc}(6a)$, $\mathrm{Loc}(8)$ qu'il invoque, ni ne
définit $W_{2}$ et $W_{\infty}$. Un lecteur qui voudrait vérifier les
raisonnements de la section 7 devra les prendre ailleurs.

\end{document}
